% BORRADOR de la Introducción (no es c1.tex). Decisiones del 24-oct-2026; P6 y P7 acordados con el usuario.

% --- P1 gancho (119 palabras)
Generating images that are not in a dataset but could belong to it is a way to overcome the scarcity of data that limits many applications: a model that learns the essential features of a small set of images can produce new samples of them \cite{Hu2025}. For example, in agriculture, images generated with a diffusion model improved the detection of vineyard shoots by up to 28.65\% in mean average precision \cite{Hirahara2025}, and generated UAV images improved the detection and counting of sorghum panicles \cite{Cai2022}; and in medical imaging, a classifier trained only with synthetic brain MRI slices reached an AUC of 0.779, against 0.869 with real data, a difference that could not be distinguished at that sample size \cite{Fallah2026}.

% --- P2 puente, Costa Rica y problema (corregido 24-oct) (157 palabras)
This scarcity is especially marked in agriculture, where the images that can be collected are few and homogeneous. Pineapple is a clear case: in Costa Rica it is grown on more than 40\,000 hectares, with 60\,000 to 100\,000 plants per hectare \cite{Castro2019}, and its monitoring for signs of disease or nutritional deficiency remains heavily reliant on slow, manual labor \cite{Gonzalez2026}. Yet a drone campaign covers a few flights over a single field, with one type of soil and similar illumination \cite{Gonzalez2026}; a detector trained in one domain does not necessarily work in another \cite{Bumbaca2025}; and annotating the images is slow and costly \cite{Bah2018,Cai2022}. Synthetic images could enlarge such datasets without new flights, with more instances of the scenes that the generator has learned, and tasks such as counting plants from drone images \cite{Kitano2019,Shiu2023} could benefit from a richer training set. Whether they add real variety, however, depends on the generator and on the latent space it works in.

% --- P3 por que el VAE (reordenado 24-oct) (213 palabras)
The use of latent diffusion models has made it far cheaper to generate realistic images, and behind that efficiency there is a component that rarely gets attention: the autoencoder that compresses each image into a latent space \cite{Rombach2022}. The diffusion model works in that compressed space, so the autoencoder is the foundation of the whole process: what it cannot represent cannot be recovered during diffusion. In practice, a pretrained autoencoder is the standard strategy \cite{Zheng2025}, but with few data from a specific domain it is not enough: the reconstructions of a VAE pretrained for latent diffusion had a precision of 0.05 on pineapple field imagery, which led its authors to conclude that it generalizes poorly and needs to be fine-tuned \cite{XieLi2025VAE}. The autoencoder therefore has to be adapted with the few images available, and its design becomes a decision with consequences: in that regime, a continuous VAE preserved the structure of a pineapple field where a discrete one failed, although the fine detail tends to be lost \cite{Gonzalez2026}. Moreover, that decision cannot be guided by reconstruction alone, because a better reconstruction does not guarantee a better generation \cite{Yao2025,Zhong2026}. The VAE is also a standard and controllable component: the data it sees, its losses, and its structure can be changed and their consequences measured.

% --- P4 que existe (144 palabras)
In the work closest to this one, the architectural adaptations (a continuous latent space and a convolutional U-Net without attention) were derived manually through systematic experimentation, a process its authors identify as a candidate for automated search \cite{Gonzalez2026}. Beyond it, the design of the autoencoder for latent diffusion has been studied one aspect at a time and, mostly, with large datasets: the number of latent channels \cite{Yao2025,Sehwag2025}, the high-frequency content of the latent space \cite{Skorokhodov2025}, masked pretraining of the autoencoder \cite{Chen2025,Lee2025LDMAE}, and the structure of the latent space of pretrained autoencoders in medical imaging \cite{Dombrowski2026}. With small datasets, the works found varied a single choice: a continuous against a discrete latent space in pineapple imagery \cite{Gonzalez2026}, the preprocessing of a small medical cohort \cite{Fallah2026}, or the weight of the regularization and the size of the codebook, evaluated by reconstruction and not by generation \cite{XieLi2025VAE}.

% --- P5 brecha y evidencia (166 palabras)
According to the literature reviewed (Chapter~II, Section~\ref{sec:litsearch}), there are no statistically sound comparisons of how several treatments of the training data and modifications of the structure of the VAE change the quality of the images generated with latent diffusion when the dataset is small, homogeneous, and multimodal. Preliminary experiments (Chapter~V) show why they are needed: the FID of two checkpoints of the same training differed by up to nine points, as much as the differences observed between configurations, and a modification that improved the reconstruction, data augmentation, did not improve the generation. Without such comparisons, the design proceeds by trial and error, and each trial requires training an autoencoder and a generator, which consumes computing resources; an apparent improvement may be due to the variability of training and not to the modification, and what is optimized may not be what matters. It also remains open whether generators trained on few and homogeneous images produce real diversity or only reproduce the restricted variability of the original data.

% --- P6 contribucion (decidido 24-oct) (116 palabras)
This work addresses that gap. Its contribution is to characterize, with replicated experiments, the effect that treatments of the training data of a variational autoencoder (augmentation, random masking, selection of images, and modality) and modifications of its structure (a variational output layer and an adversarial loss) have on the quality of the aerial pineapple images generated with DDPM and Flow Matching, estimating its direction, size, and interactions, and identifying which effects are favorable, which are harmful, and which are negligible. For a reader who wants to generate synthetic images from a small dataset, the result indicates which modifications of the data and of the VAE are worth making, and which are not, within the conditions studied.

% --- P7 enfoque y alcance (123 palabras)
The study follows a factorial design in stages: a screening of the factors with several independent seeds, which separates the effect of each modification from the variability of training, a refinement of those that matter, and a confirmation with Flow Matching (Chapter~IV). It is limited to one dataset of pineapple fields: the transfer to other crops or to other domains, such as medical imaging, the systematic comparison with other types of latent spaces (explored only in a feasibility study), and the explanation of the mechanisms behind the effects are left as future work, together with labeling the generated images and using them to train a counting model. The work continues the research line of the group on pineapple imagery captured by drones \cite{XieLi2024,Gonzalez2026}.

% --- P8 mapa (59 palabras)
The rest of this document is organized as follows. Chapter~II presents the application context, the theoretical framework, and the related work, including the literature search. Chapter~III states the problem, the hypothesis, the objectives, the research questions, and the scope. Chapter~IV describes the methodology and the design of the experiments. Chapter~V reports the initial results, and Chapter~VI presents the workplan.

